\documentclass[10pt,letterpaper,fontset=fandol]{ctexart}

% ---------- packages ----------
% NOTE: 用 XeLaTeX 编译。不要加 inputenc / fontenc —— ctex 自己处理 CJK。
\usepackage{microtype}
\usepackage{amsmath,amssymb}
\usepackage{graphicx}
\usepackage{booktabs}
\usepackage{multirow}
\usepackage{xcolor}
\usepackage{natbib}
\usepackage{hyperref}
\usepackage[capitalise]{cleveref}
\usepackage{enumitem}
\usepackage[margin=1in]{geometry}
\usepackage{xspace}

\hypersetup{
    colorlinks=true,
    linkcolor=blue!60!black,
    citecolor=blue!60!black,
    urlcolor=blue!60!black
}

\newcommand{\todo}[1]{\textcolor{red}{\textbf{[待办：#1]}}}
\newcommand{\method}{\textsc{hippoVLA}\xspace}
\newcommand{\sTMB}{\textsc{sTMB}\xspace}

% ---------- title ----------
\title{\method：面向视觉--语言--动作模型的\\
类海马体时空解耦框架}

\author{
    匿名作者\\
    \texttt{\{anonymous\}@institution}
}

\date{}

\begin{document}
\maketitle

% =====================================================================
\begin{abstract}
视觉--语言--动作（VLA）模型继承了视觉--语言模型（VLM）的单一 Transformer 主干，因而需要在同一个上下文窗口内同时承担两类在认知上截然不同的功能：单帧的\emph{空间定位}与跨帧的\emph{时间整合}。我们认为，这种纠缠是当前 VLA 长程执行能力薄弱、空间推理不稳定这两类失败的结构性根源——两类功能在有限的表征容量上相互争夺。
受人类记忆的互补学习系统（Complementary Learning Systems, CLS）理论启发——新皮层缓慢地巩固统计规律，而海马体则快速编码情景上下文——我们提出 \method，一种将上述两条流式处理路径显式解耦的类脑 VLA 框架。\method 由一个空间增强的 VLM 主干（对应背侧通路皮层）与一个类海马体的短期记忆模块（Short-Term Memory Bank，\sTMB）组成。\sTMB 包含三项关键设计：（i）模式分离的可学习记忆槽；（ii）稀疏的\emph{theta 节律}时间采样；（iii）一种新颖的\emph{空间显著性感知门控}，仅将空间相关的 token 选择性写入记忆。在 LIBERO 4-in-1 上，\method 在 Object/Spatial/Goal/Long 四个子集上分别取得 $99.8/94.4/96.4/96.0$；在 CALVIN ABC$\rightarrow$D 长程任务排行榜上获得平均完成长度 $4.38$（排名第 3），相较同主干无记忆基线绝对提升 $+0.68$。消融实验表明：空间主干与记忆模块单独使用时均不充分，二者\emph{协同}才能突破各自的能力上限——这恰好对应 CLS 框架中两类记忆系统的分工。
\end{abstract}

% =====================================================================
\section{引言}
\label{sec:intro}

视觉--语言--动作（VLA）模型通过在多模态 Transformer 之上挂接动作解码头，将视觉--语言模型（VLM）拓展至具身控制~\citep{black2024pi0,kim2024openvla,brohan2023rt2}。这一范式在单步操作任务上取得了显著成功，却在两种困难场景下持续失利：需要跨多帧整合信息的\emph{长程任务}，以及需要精确三维定位的\emph{空间敏感任务}。我们指出一个尚未充分探讨的结构性原因：同一组自注意力堆叠被同时要求完成两类功能，而这两类功能在固定的上下文窗口内此消彼长。

\paragraph{单一通路中的空间--时间冲突。}
帧级 VLA token 同时承担两项使命：将语言定位到当前场景，以及作为唯一的载体把过去的视觉历史传递到未来。随着上下文增长，注意力会稀释空间细节；而当模型为空间精度而调优时，长距离时间模式则更难被保留。我们在 LIBERO 与 CALVIN 上同时观察到这两类失效：纯粹的强空间主干在长程任务上止步于 $90\%$ 以下（见 \Cref{tab:libero-cross}），而专为长程任务优化的方法又在空间任务上让步。

\paragraph{一种类脑替代方案。}
生物智能体并不靠单一同质网络解决这一冲突。互补学习系统（CLS）假说~\citep{mcclelland1995cls,kumaran2016cls}认为存在两类相互作用的记忆系统：一类是\emph{慢}的新皮层系统，跨众多经验学习可泛化的表征；另一类是\emph{快}的海马体系统，快速编码独立情景，并在巩固与检索阶段进行回放。海马体并非被动缓冲：齿状回的模式分离产生稀疏、去相关的记忆痕迹；theta--gamma 耦合为编码施加离散时间节律；前额叶的反馈则决定哪些进入的事件值得托付给记忆~\citep{rolls2013mechanisms}。

\paragraph{本工作。}
我们提出 \method，一种与上述分工相对应的 VLA 架构。空间增强主干（RynnBrain-8B）扮演背侧通路 / 新皮层的角色，提供丰富的单帧空间定位；新提出的\textbf{短期记忆模块}（\sTMB）扮演海马体角色，维护一组固定大小、模式分离的记忆槽，以稀疏时间节律跨帧吸收信息。关键之处在于，\sTMB 的写入由一个\emph{空间显著性感知门控}决定，该门控查询空间主干，决定哪些 token 值得记忆。这使得两条通路是\emph{耦合}而非简单拼接的：空间系统决定了时间系统应当存储什么。

\paragraph{贡献。}
\begin{itemize}[leftmargin=*]
    \item 我们诊断出单一 VLA 主干中存在结构性的空间--时间纠缠问题，并提出受 CLS 启发的解耦设计。
    \item 我们提出 \sTMB，一种类海马体的记忆模块，包含模式分离的记忆槽、稀疏时间采样以及创新的空间显著性感知写入门控。
    \item \method 在 LIBERO 4-in-1 四个子集上取得 $99.8/94.4/96.4/96.0$；在 CALVIN ABC$\rightarrow$D 上达到平均完成长度 $4.38$（排名第 3），相较同主干无记忆基线绝对提升 $+0.68$。
    \item 消融实验表明空间主干与记忆模块\emph{各自不充分但合为充分}，证实了 CLS 解读：两条通路是协同而非冗余。
\end{itemize}

% =====================================================================
\section{相关工作}
\label{sec:related}

\paragraph{视觉--语言--动作模型。}
近年的 VLA 系统——OpenVLA~\citep{kim2024openvla}、$\pi_0$~\citep{black2024pi0}、RT-2~\citep{brohan2023rt2}、GR00T~\citep{nvidia2025gr00t}——在预训练 VLM 之上挂接动作解码器（离散 token、MLP 回归或 diffusion）。绝大多数将 VLM 视作黑盒特征提取器，未在原生上下文窗口之外显式建模时间上下文。

\paragraph{机器人学习中的记忆。}
循环或 Transformer 策略通过隐藏状态或上下文实现隐式记忆~\citep{brohan2023rt2,reed2022gato}。显式记忆的探索包括外部检索~\citep{wu2022memorizing}、槽注意力~\citep{locatello2020slot}与记忆 token~\citep{bulatov2022rmt}，但很少有方法配合显式门控、并将该门控的判定下放给一条独立的空间通路。同期工作如 \todo{补充近期记忆增强 VLA 引用}虽扩展了上下文窗口，但并未在容量层面上解耦空间与时间。

\paragraph{类脑人工智能。}
深度学习中的海马体模型包括可微神经计算机~\citep{graves2016dnc}、神经情景控制~\citep{pritzel2017nec}与 CLS 风格的回放~\citep{kumaran2016cls,parisi2019cls}。本工作的差异主要体现在两点：其一，我们以\emph{架构}方式而非\emph{流程}方式（如回放缓冲）实例化 CLS——两条平行的通路；其二，门控由姊妹网络条件化，而非自条件化。

\paragraph{空间基础模型。}
诸如 SpatialBot、RynnBrain 等空间增强 VLM 通过点 / 深度对齐的数据训练，提升几何推理能力。我们采用 RynnBrain-8B 作为背侧通路组件，但关键并不在于"使用"——而在于将它与一个能够消费其空间显著性信号的记忆模块\emph{耦合}。

% =====================================================================
\section{方法}
\label{sec:method}

\method 由三部分组成（\Cref{fig:arch}）：空间通路（\Cref{sec:spatial}）、类海马体的 \sTMB（\Cref{sec:stmb}）以及动作解码器。下面分别介绍，并阐述三者的整合方式。

\begin{figure}[t]
\centering
\todo{图：双通路架构示意。左：空间通路（RynnBrain-8B）输出每帧的视觉 token 与一张空间显著性图。右：含 $N$ 个槽的 \sTMB，theta 节律采样的写路径，空间显著性感知门控。底部：动作解码器同时读取两条通路。}
\caption{\method 架构。空间通路（左）输出单帧 token 与显著性信号；\sTMB（右）维护一组固定数量、模式分离的记忆槽，以稀疏时间节律更新；其写门由空间显著性条件化，因此只有空间相关的 token 才会进入记忆。动作解码器同时读取两条通路，对应 CLS 中海马体与新皮层之间的协同。}
\label{fig:arch}
\end{figure}

\subsection{认知建模视角}
\label{sec:framing}
我们将 $t$ 时刻的 VLA 前向过程写作 $a_t = \pi(o_t, h_t)$，其中 $o_t$ 是当前观测，$h_t$ 概括历史。在单一 Transformer 中，$h_t$ 隐含于上下文窗口内仍存在的过去 token，并且与计算 $o_t$ 空间特征的同一组参数共享。沿用 CLS 思路，我们将 $\pi$ 显式分解为快速情景通路 $\mathcal{H}$（海马体）与慢速语义通路 $\mathcal{N}$（新皮层）：
\begin{equation}
    a_t = D\bigl(\,\mathcal{N}(o_t),\; \mathcal{H}(o_{\le t})\,\bigr),
\end{equation}
其中 $D$ 为轻量动作解码器。两条通路是\emph{耦合}的：$\mathcal{H}$ 写入记忆需经 $\mathcal{N}$ 同意，借由空间显著性感知门控（\Cref{sec:gate}）实现。

\subsection{空间通路：RynnBrain-8B}
\label{sec:spatial}
我们采用 RynnBrain-8B 作为空间主干。对每一帧 $o_t$，它输出与语言 token 对齐的视觉 token $V_t \in \mathbb{R}^{N_v \times d}$。RynnBrain 通过点 / 深度对齐的监督完成预训练，因此 $V_t$ 内含隐式的三维先验；我们额外抽取每个 token 的空间显著性 $s_t \in [0,1]^{N_v}$，其定义为 RynnBrain 空间定位通路与一条非空间基线通路之间残差的 L2 范数 \todo{最终确定定义；候选方案：对 \textsc{ChainOfPoint} 的 [PT] token 做 attention rollout}。$s_t$ 是从 $\mathcal{N}$ 流向 $\mathcal{H}$ 的唯一信号。

\subsection{\sTMB：类海马体情景记忆}
\label{sec:stmb}

\paragraph{记忆库即模式分离的槽。}
\sTMB 维护固定尺寸的记忆 $M \in \mathbb{R}^{B \times K \times d}$，含 $K=64$ 个可学习槽嵌入，初始化为独立同分布高斯，并在同一情景内跨帧持续更新。槽数远小于每帧 token 数（$K \ll N_v$），形成一个类似齿状回模式分离的瓶颈：模型必须把情景内容压缩到不同的槽中，而非存储原始帧。

\paragraph{theta 节律稀疏采样。}
海马体编码由 theta--gamma 耦合调控：仅一部分进入的感官帧才会被托付给长期记忆。我们用一个可调的\emph{采样间隔} $\tau$ 模拟该机制：\sTMB 每 $\tau$ 帧才执行一次写入，而非每帧都写。在实验中（\Cref{sec:exp-ablation}）$\tau=10$ 优于 $\tau=5$，与"更稀疏、更去相关的样本构成更佳情景锚点"的观点一致。

\paragraph{含空间显著性感知门控的写路径。}
\label{sec:gate}
在写入步 $t$，记忆槽通过交叉注意力关注进入的视觉 token：
\begin{equation}
    \tilde M_t = \mathrm{CrossAttn}(M_{t-1},\, V_t + e_t),
\end{equation}
其中 $e_t$ 为可学习的时间位置嵌入。标准的门控更新为 $g = \sigma(W_g[M_{t-1};\tilde M_t])$，并按 $M_t = g \odot \tilde M_t + (1{-}g)\odot M_{t-1}$ 融合。我们在此基础上引入显著性条件化的乘法项：
\begin{equation}
    g_t \;=\; \sigma\!\bigl(W_g\, [M_{t-1};\,\tilde M_t]\bigr) \;\odot\; \phi(\bar s_t),\qquad
    M_t = g_t \odot \tilde M_t + (1 - g_t) \odot M_{t-1},
    \label{eq:gate}
\end{equation}
其中 $\bar s_t \in [0,1]$ 为 $V_t$ 的均值空间显著性，$\phi$ 为一个返回逐槽标量的小型 MLP。该设计的直觉直接：当空间通路告知当前帧空间信息贫瘠（如相机运动、遮挡、无可操作物体），门控关闭、\sTMB 保持原状；当帧空间显著（物体新近可见、夹爪接近接触）时，门控打开。这正是 CLS 中前额叶--海马体门控的架构对应。

\paragraph{读路径。}
预测 $t$ 时刻动作时，动作解码器的查询 token $Q_t$ 通过交叉注意力检索 \sTMB：
\begin{equation}
    \hat Q_t = Q_t + \mathrm{CrossAttn}(Q_t, M_t).
\end{equation}
解码器因此消费一个同时包含\emph{当前空间定位}（来自 $\mathcal{N}$）与\emph{情景上下文}（来自 $\mathcal{H}$）的融合表征。

\subsection{训练与推理}
\sTMB 与动作头联合训练；空间主干以较小学习率微调。$M$ 在每个情景开始时由可学习原型初始化，并在情景之间重置。\sTMB 在主干之上新增约 \todo{X}M 参数，占总参数的 $<$\todo{X}\%。

% =====================================================================
\section{实验}
\label{sec:exp}

我们在两个互补的具身基准上评估 \method：\textbf{LIBERO 4-in-1}~\citep{liu2023libero}（在 Spatial / Object / Goal / Long 四个子集上联合训练单一策略）与 \textbf{CALVIN ABC$\rightarrow$D}~\citep{mees2022calvin}（语言条件下的长程操作，主要指标为五条指令中平均完成数）。

\subsection{实验设置}
\paragraph{训练。} 在 LIBERO 4-in-1 上联合训练，batch size 16，优化步数 20k--60k，AdamW。主干：用于记忆消融的 starVLA-Qwen3VL-4B；用于跨主干对比与 CALVIN 的 RynnBrain-8B-OFT。\sTMB 使用 $K=64$ 槽，$\tau \in \{5,10\}$，记忆长度 5。
\paragraph{推理。} \sTMB 在情景开始时重置；记忆在整个情景期间持续。

\subsection{LIBERO 4-in-1：记忆消融}
\label{sec:exp-libero-ablation}
\Cref{tab:libero-mem} 汇报相同 Qwen3VL-4B 主干下不同记忆配置的成功率。$\tau=10$ 的 \sTMB 相较无记忆基线在 Long 上提升 $+3.0$、Spatial 上提升 $+1.2$，并以 20k 步达到与 50k 步官方权重相当的水平。$\tau=5$ 与 $\tau=10$ 的对比与 theta 节律解读一致：更稀疏、更去相关的写入产生更优的情景记忆。

\begin{table}[t]
\centering
\small
\caption{LIBERO 4-in-1 记忆消融（成功率 \%）。同一 Qwen3VL-4B 主干，仅记忆配置变化；$\tau=10$ 的 \sTMB 在四个子集中的三个上取得最佳。}
\label{tab:libero-mem}
\begin{tabular}{lcccccccc}
\toprule
设置 & BS & 步数 & 主干 & 模块 & Object & Spatial & Goal & Long \\
\midrule
官方                             & --- & 50k & Qwen3VL-4B & ---       & 99.6 & ---  & 99.0 & 98.6 \\
本机复现                          & --- & 50k & Qwen3VL-4B & ---       & 99.6 & 93.2 & 98.6 & 93.0 \\
+ \sTMB ($\tau{=}5$)             & 16  & 20k & Qwen3VL-4B & len 5 & 99.6 & 94.2 & 98.2 & 95.8 \\
+ \sTMB ($\tau{=}10$) \textbf{[本工作]} & 16  & 20k & Qwen3VL-4B & len 5 & \textbf{99.8} & \textbf{94.4} & 96.4 & \textbf{96.0} \\
\bottomrule
\end{tabular}
\end{table}

\subsection{LIBERO：跨主干对比}
\label{sec:exp-libero-cross}
\Cref{tab:libero-cross} 表明，在 RynnBrain-8B-OFT 主干之上引入 \sTMB 后 Spatial 提升 $+1.0$，同时保持空间主干在 Object/Goal/Long 上的强表现。这正是 CLS 假说在实证上的核心特征：单独的空间通路即便强大也在 Spatial 上受限，而单独的记忆通路（搭配较弱主干）则受限于绝对上限——只有两者并用才能在两个轴上同时饱和。

\begin{table}[t]
\centering
\small
\caption{LIBERO 跨主干对比。\sTMB 把空间能力强的 RynnBrain 主干推过其无记忆上限。}
\label{tab:libero-cross}
\begin{tabular}{lccccc}
\toprule
方法 & 备注 & Spatial & Object & Goal & Long \\
\midrule
Qwen2.5-VL-GR00T-LIBERO-4in1            & 本机复现     & 98.4 & 99.0 & 96.2 & 91.0 \\
RynnBrain-8B-OFT（无记忆）              & step 60k    & 94.4 & \textbf{100.0} & \textbf{98.0} & \textbf{95.8} \\
\textbf{\method}（RynnBrain + \sTMB）   & step 60k    & \textbf{95.4} & 99.6 & 97.2 & 93.8 \\
\bottomrule
\end{tabular}
\end{table}

\subsection{CALVIN ABC$\rightarrow$D 长程任务}
\label{sec:exp-calvin}
\Cref{tab:calvin} 给出完整的 CALVIN 长程任务排行榜。\method（RynnBrain + \sTMB）取得平均完成长度 $4.38$，在已发表方法中排名第 3——更关键的是，相较同主干的无记忆消融（3.70）绝对提升 $+0.68$。这是我们整体流水线中单个组件的最大贡献。长程执行恰好是单一 Transformer 上下文耗尽的场景；\sTMB 通过在语言上下文之外维护情景状态而完全绕开了这一瓶颈。

\begin{table}[t]
\centering
\small
\caption{CALVIN ABC$\rightarrow$D 长程任务排行榜。Avg.\ Len.\ 为连续完成的指令数（最大 5）的均值。\method 排名第 3，比自身的无记忆消融高 $+0.68$。}
\label{tab:calvin}
\begin{tabular}{rlcccccc c}
\toprule
\# & 方法 & Step 1 & Step 2 & Step 3 & Step 4 & Step 5 & Avg.\ Len. \\
\midrule
1 & FLOWER                          & 99.4 & 95.8 & 90.7 & 84.9 & 77.8 & 4.53 \\
2 & UniVLA                          & 98.9 & 94.8 & 89.0 & 82.8 & 75.1 & 4.41 \\
3 & \textbf{\method（本工作）}      & 99.1 & 94.4 & 88.4 & 81.6 & 74.1 & \textbf{4.38} \\
4 & SeeR-Large                      & 96.3 & 91.6 & 86.1 & 80.3 & 74.0 & 4.28 \\
5 & GR-MG                           & 96.8 & 89.3 & 81.5 & 72.7 & 64.4 & 4.04 \\
6 & MoDE                            & 96.2 & 88.9 & 81.1 & 71.8 & 63.5 & 4.01 \\
7 & RynnBrain（无记忆，本工作）      & 88.7 & 79.9 & 72.7 & 67.3 & 61.0 & 3.70 \\
8 & GHIL-Glue                       & 95.2 & 88.5 & 73.2 & 62.5 & 49.8 & 3.69 \\
9 & RoboUniView                     & 94.2 & 84.2 & 73.4 & 62.2 & 50.7 & 3.64 \\
10 & Diffusion Transformer Policy   & 94.5 & 82.5 & 72.8 & 61.3 & 50.0 & 3.61 \\
11 & CLOVER                         & 96.0 & 83.5 & 70.8 & 57.5 & 45.4 & 3.53 \\
12 & 3D Diffuser Actor              & 92.2 & 78.7 & 63.9 & 51.2 & 41.2 & 3.27 \\
13 & GR-1                           & 85.4 & 71.2 & 59.6 & 49.7 & 40.1 & 3.06 \\
14 & DeeR                           & 86.2 & 70.1 & 51.8 & 41.5 & 30.4 & 2.82 \\
15 & SuSIE                          & 87.0 & 69.0 & 49.0 & 38.0 & 26.0 & 2.69 \\
16 & RoboFlamingo                   & 82.4 & 61.9 & 46.6 & 33.1 & 23.5 & 2.47 \\
17 & SPIL                           & 74.2 & 46.3 & 27.6 & 14.7 &  8.0 & 1.71 \\
18 & HULC                           & 41.8 & 16.5 &  5.7 &  1.9 &  1.1 & 0.67 \\
\bottomrule
\end{tabular}
\end{table}

\subsection{消融实验}
\label{sec:exp-ablation}

\paragraph{(a) 通路单独使用 vs.\ 联合使用。} 我们对比四种配置：(i) 仅主干；(ii) 主干 + 不带显著性门的 \sTMB（即标准 GRU 风格门）；(iii) 主干 + 带显著性门的 \sTMB；(iv) 显著性\emph{跨帧打乱}作为对照。\todo{运行并填表；预期 (iii) > (ii) > (i)，且 (iv) 退化至 (i)，从而证明显著性信号承载真实内容。}

\paragraph{(b) 采样间隔 $\tau$。} \Cref{tab:libero-mem} 已显示 $\tau=10 > \tau=5$。\todo{补 $\tau \in \{1, 3, 5, 10, 20\}$ 的扫描以画出曲线。}

\paragraph{(c) 槽数 $K$。} \todo{$K \in \{8, 16, 32, 64, 128\}$ 消融；预期为带峰的圆顶曲线，峰值约在 $K=64$（模式分离论点预测最佳点）。}

\paragraph{(d) 槽探针分析。} \todo{在 rollout 后探针单独的记忆槽，证明部分槽专门编码物体身份、另一部分槽专门编码夹爪--物体的空间关系，对应位置细胞与物体细胞的分工。}

\subsection{定性分析}
\Cref{fig:rollouts} 展示了四个 LIBERO 子集上的代表性 rollout。\todo{补充图：四组帧序列；高亮无记忆基线在遮挡后失去物体身份、而 \method 借助 \sTMB 恢复的案例。}

\begin{figure}[t]
\centering
\todo{图：四条 rollout 帧条（Spatial / Object / Goal / Long），每帧下方标注 \sTMB 写门是开启还是关闭。}
\caption{\method 在 LIBERO 上的定性 rollout。下方条形指示空间显著性门控状态：开启（实心）表示当前帧被写入 \sTMB，关闭（斜线）表示未写入。}
\label{fig:rollouts}
\end{figure}

% =====================================================================
\section{讨论}
\label{sec:discussion}

\paragraph{为什么显著性门控起作用？}
原则上，单纯由数据驱动的门控（在 $[M;\tilde M]$ 上做 sigmoid）也能学到同样的行为，但它必须只依靠稀疏的动作预测梯度学得。条件化于 $s_t$ 注入了一种归纳偏置：``信任空间通路对何为值得记忆的判断''。这正是哺乳动物记忆中归属于前额叶--海马体反馈的角色。

\paragraph{局限。}
\sTMB 仅承担短期记忆，并未跨情景巩固。CLS 还预测了一个更慢的过程，将情景内容蒸馏入皮层；我们将这部分工作留作后续研究（例如基于回放的主干微调）。当前显著性的定义是启发式的；与动作预测联合训练的可学习显著性头是一个自然的延展。

\paragraph{关于复用预训练空间主干。}
一个自然的疑问是 \method 是否仅是对既有空间主干（RynnBrain）的二次封装。\Cref{tab:libero-cross} 与 \Cref{tab:calvin} 中的无记忆消融表明并非如此：单独的主干即便经过微调也存在上限；增益来自空间通路与 \sTMB 之间通过显著性门控的\emph{耦合}。

% =====================================================================
\section{结论}

我们提出了 \method——一种受海马体启发的视觉--语言--动作模型，借助 CLS 风格的双通路架构将空间定位与时间整合显式解耦。所提出的短期记忆模块通过模式分离的槽、稀疏 theta 节律采样以及空间显著性感知门控，将一个强空间主干推过其无记忆上限，在 LIBERO 与 CALVIN 上均取得显著提升。我们将这视为一个重要的方向：让 VLA 系统的\emph{架构}（而非仅仅是训练数据）反映其所要模拟的具身智能体的认知结构。

% =====================================================================
\small
\bibliographystyle{plainnat}
\begin{thebibliography}{99}

\bibitem[Black et al.(2024)]{black2024pi0}
K.~Black et al.\ \newblock $\pi_0$: A vision-language-action flow model for general robot control. 2024.

\bibitem[Brohan et al.(2023)]{brohan2023rt2}
A.~Brohan et al.\ \newblock RT-2: Vision-language-action models transfer web knowledge to robotic control. 2023.

\bibitem[Bulatov et al.(2022)]{bulatov2022rmt}
A.~Bulatov et al.\ \newblock Recurrent memory transformer. \emph{NeurIPS}, 2022.

\bibitem[Graves et al.(2016)]{graves2016dnc}
A.~Graves et al.\ \newblock Hybrid computing using a neural network with dynamic external memory. \emph{Nature}, 2016.

\bibitem[Kim et al.(2024)]{kim2024openvla}
M.~Kim et al.\ \newblock OpenVLA: An open-source vision-language-action model. 2024.

\bibitem[Kumaran et al.(2016)]{kumaran2016cls}
D.~Kumaran, D.~Hassabis, J.~L.~McClelland. \newblock What learning systems do intelligent agents need? Complementary learning systems theory updated. \emph{Trends in Cognitive Sciences}, 2016.

\bibitem[Liu et al.(2023)]{liu2023libero}
B.~Liu et al.\ \newblock LIBERO: Benchmarking knowledge transfer for lifelong robot learning. 2023.

\bibitem[Locatello et al.(2020)]{locatello2020slot}
F.~Locatello et al.\ \newblock Object-centric learning with slot attention. \emph{NeurIPS}, 2020.

\bibitem[McClelland et al.(1995)]{mcclelland1995cls}
J.~L.~McClelland, B.~L.~McNaughton, R.~C.~O'Reilly. \newblock Why there are complementary learning systems in the hippocampus and neocortex. \emph{Psychological Review}, 1995.

\bibitem[Mees et al.(2022)]{mees2022calvin}
O.~Mees et al.\ \newblock CALVIN: A benchmark for language-conditioned policy learning for long-horizon robot manipulation tasks. \emph{IEEE RA-L}, 2022.

\bibitem[NVIDIA(2025)]{nvidia2025gr00t}
NVIDIA. \newblock GR00T N1: An open foundation model for generalist humanoid robots. 2025.

\bibitem[Parisi et al.(2019)]{parisi2019cls}
G.~I.~Parisi et al.\ \newblock Continual lifelong learning with neural networks: A review. \emph{Neural Networks}, 2019.

\bibitem[Pritzel et al.(2017)]{pritzel2017nec}
A.~Pritzel et al.\ \newblock Neural episodic control. \emph{ICML}, 2017.

\bibitem[Reed et al.(2022)]{reed2022gato}
S.~Reed et al.\ \newblock A generalist agent. \emph{TMLR}, 2022.

\bibitem[Rolls(2013)]{rolls2013mechanisms}
E.~T.~Rolls. \newblock The mechanisms for pattern completion and pattern separation in the hippocampus. \emph{Frontiers in Systems Neuroscience}, 2013.

\bibitem[Wu et al.(2022)]{wu2022memorizing}
Y.~Wu et al.\ \newblock Memorizing transformers. \emph{ICLR}, 2022.

\end{thebibliography}

\end{document}
